Six out of seven G7 countries include standardized exams for university admission. 

This study documents the reordering of student's accessibility to elite university programs through a switch from standardized exam to non-standardized testing. I study this question by documenting grade inflation in the UK during the COVID-19 pandemic, caused by the cancellation of their standardized exam (A levels). Two months before the exam date, the British government banned in-person exams and allowed students to use their teachers predicted grade as their final A-level grade. Without any unified procedure for assessing grades, teachers provided grades using their personal knowledge and experience with the students. The resulting grade distribution moved dramatically upward, with the number of top grades increasing two times. The grade shift was mainly driven by grade inflation, as students had mostly finished their educational curriculum at the time of the cancellation, and students had already agreed with universities on the grade requirements for securing placements. 

In the first part of this paper, I examine how the grading method change changed the probability of obtaining top grades across the students backgrounds. Namely, I focus on previous academic achievements, demographic attributes (race, gender, and parental income), subject content, and school of attendance. Achievements in final exams are strongly connected with background of students as backgrounds motivate and prepare students for their exams in different ways. The gap is persistent across demographic groups, despite standardized exams not taking into account the background of the students. I investigate the full extent of how replacing standardized exams with non-standardized tests benefits students differently.  Firstly, I investigate whether the students were graded in proportion to their previous academic performance before sitting their final exams. The British higher education system mandates all students to take a standardized test before starting their high-school education. I test whether the extent of the increase in obtaining higher grades was higher for students with academic background. Furthermore, I distinguish between resorting of students within classrooms and the overall student distribution.  

Secondly, I investigate the demographic attributes related to higher increase in obtaining top grades. As teachers may be biased towards certain attributes of students, such biases may directly reflect into the differences in achievement in their final exams. I focus on gender, race as teachers can easily distinguish between demographic groups. I also investigate whether parental occupation reflects into differences, as parents influence the final decision making of teachers. Thirdly, I investigate whether the content of the final exams created differences in the students likelihood of obtaining a top grade. Quantitative subjects are easier to predict the final output of teachers are quantitatively informed about students scores in previous internal tests, whereas non-quantitative subject may be more difficult to grade given the subjectiveness of their content. However, quantitative subjects are more often involved in ambitious applications, as top university programs often require applicants to submit grades in competitive subjects.  

Finally, I investigating the difference in the extent of increasing students across the full set of secondary schools. As privately funded schools attract pupils by placing their students in elite university programs, they may have larger incentive to over-report their students grades at the expense of academic integrity. I test whether the selective feature of private schools explain the extent of their inflation patterns by comparing the private schools inflation patterns to selective public schools in the UK (Grammar schools). I also investigate how the extent of the grade increase was associated with the quality of the school. Lower quality schools may provide higher grades to students to compensate for their lack of educational resources. My analysis documents the beneficiaries of the grading scheme change by decomposing the effect on students, subjects, and schools.

In the second part of the paper, I document the responses of universities by adjusting their admission thresholds. British universities send admission offers with grade cut-offs that specify the grades that students must meet to secure placement at their institution. Universities write the offer letters to each universities using before students sits their exams. The offer letters constrain the universities to legally oblige to accept students passing the grade threshold. Since tuition caps constrain British universities ability to control admission volume, admission thresholds dominated universities means to respond to the sudden changes in the overall grade distribution. Universities may trade off quality of students with filling their capacity to accept students. Previous literature documents how students care about the quality of their incoming cohorts, and universities set admission thresholds to maintain high quality of students while filling their teaching capacity. 

In the final part of the paper, I examine how the altered grade distribution changed the composition of students at universities. Despite the accusations that private schools inflated grades more than public schools, the school's impacts on compositions on students in university programs might be limited due to the students' already high chances of securing a placement in an elite program. On the other hand, students from public schools increase their representation in elite programs due to the large volume of students on the margins of receiving a high grade. Furthermore, the different extent of grade inflation across student groups and subject choice within state school may also reflect into the student composition at universities. I also take into account application patterns of different student groups, as previous literature documents underrepresented students apply much cautiously than other students.

\

I estimate grade inflation for each student using administrative data on entry levels A-level. The British higher education system is well suited for measuring grade inflation, as all students take standardized exams (GCSE) when they start their 16-18 education. The standardized measure of students academic performance allows me to compare the increments in their grade on a common measure of academic ability for the entire population of higher education applicants in the UK. Unsurprisingly, the GCSE scores accurately predict the results of the student grade A-levels that I test using an event-study test. I obtain the causal estimates of increments in the grade received by using an interrupted time series design that leverages the sudden announcement of the policy and the short time window for teachers to provide their predictions. Specifically, I measure grade inflation as the increments in the probability of the student receiving a top grade. Using administrative data on university application to the UK, I traces the grade outcomes of the universe of students and document the upward shift in final grade while conditioning the movements on the students' background or their previous academic performance. Namely, I estimate a high-dimensional logit model that controls for students attributes, school of attendance, and the subject of the exam. The model estimates more than 10,000 parameters, which I stabilize using shrinkage estimators. The method recovers a rich set of estimates which I use to compare the degree of grade reassignment across a vast range of student backgrounds. I also document how assignment patterns evolved within schools during the two years of the pandemic. Namely, employ a instrumental variable approach motivated by Arellano \& Bond (1998) to estimate how the degree of grade inflation in the first year affected the inflation levels in the next year. I additionally estimate the degree of inflation spillover across subject groups (Non-quantitative vs. quantitative subject groups). 

I estimate the readjustments of admission thresholds of universities by estimating the admission cutoff for each university program. British universities set a single admission cutoff rule among their entire applicant pool. Using this institutional detail, I fit a multi-dimensional kernel density function over the grade space with a binary acceptance/rejection variable to infer the admission cutoff. The approach is motivated by Porter and Yu (2015) and applied by Mountjoy (2024). I apply this method over all university program in the UK and uniquely estimate the cutoff for each university programs. Additionally, I separately estimate the cutoff across 3 periods of my data coverage; pre-pandemic, 2020, and 2021. Notably, universities were unaware of the grading method change in 2020, while in 2021 the British government did not announce the cancellation of in-person exams until the 4 months before the exam date. I document the patterns in admission threshold adjustment differed across universities. 

In the final part of my paper, I develop a quantitative framework to compare the composition of student demographic group across universities. Following the general equilibrium framework proposed by Epple, Romano, and Sieg (2006), the model is calibrated to the exact demographic distribution of the student body in the UK, and capacity constraints of each British University. The model assigns grades differently based on the grading scheme; in which the non-standardized grading scheme assignment is based on the students school of attendance, their demographics, and subject. Student choose which university to apply to based on the rank of the universities. Universities set grade thresholds to optimize profits while keeping the volume of incoming cohorts under the capacity limit. Universities differ in their marginal costs for over/under accepting students. I additionally estimate the marginal costs by regressing the degrees of threshold adjustment in 2021 on the degree of over-acceptance in 2020 for each universities. 

\

I provide evidence of grade inflation after removing the standardized exams. Substantial heterogeneity arises in the attributes of the student, the school, and the exam subject. Conditional on the academic standing of the student, the odds increase is larger for students with higher academic standings. A student with the second highest academic standing (GCSE = 8) increased their odds of receiving an A/A* by 26 percent, while students with the lowest grade (GCSE = 3) decreased their odds by 50 percent. The increase/decrease in their odds is proportional to the students' achievement scores. In probability terms, a student with good academic standing with 50\% chances of receiving an A/A increased their probability to 58 \%, while a student with worse academic standing but with 50 \% on receiving A/A* reduced their chances of receiving A/A* to 35\%. The patterns in grade assignment and student achivements is robust to including subject level controls in the regression. This suggest that the ordering of students reflected the students previous academic achievements.

I find that female and white students increased their odds of receiving high grades. Female students increased their odds of receiving A/A* by 22\% more than male students. Since the female have lower probability in obtaining high grades that males students in standardized exams, my results show that students sample White students increased their odds of receiving higher grades by 7 \% more than non-white students. I do not find evidence of 
pupils in private schools obtained top grades at [] percent, which is an increase of [] percentage than in previous years. Students attending public schools received a top grade with [] percent, a [] percentage point increase. Depending on the demographics of the students in public schools, private schools assigned the top grades by [] percentage points more than public schools. Schools in populated regions inflated grades more. A percentage increase in school numbers was associated with a inflation rate of [] percentage points higher. School inflate by assigning top grades to students taking non-quantitative or non-facilitating, a subject category recommended by elite university programs. Students taking non-quantitative and non-facilitating subjects increased their probability of receiving top grades by 18 and 17 percentage points, respectively. Quantitative / facilitative subject takers received a top grade of 5 and 8 percentage points more than their counterparts, respectively. Lastly, grade inflation occurred differently between demographic groups of students. Female students received higher grades more likely than male students by 5 percentage points, white students received higher grades more likely by 3 percentage points. Inflation occurred differently across occupational class of students, ranging from a lower bound of 11 percent points increase for the students with parents that never worked or long-term unemployed, to a upper bound of 28 percentage points for students with parents in lower supervisory or technical occupations. However, demographic-related estimates lose statistical significance once the regression includes the student's subject choice.

My quantitative model predicts university composition of incoming students changes in particular student's originating schools and average academic performance in their high school exams. Taking the acceptance threshold as constant, the share of students from private schools \emph{decreases} by [] percent. In contrast, students in public universities increase by [] percent. This reflects the underlying connection between elite programs and private school institutions, private school students occupying almost [] percent of the student body in elite programs. In contrast, the share of students from public schools in elite programs increases by [] percent. The shift of students from public schools coincides with changes in demographic composition in elite universities, as the share of white and female students increases by [] percent in elite programs. In contrast, the share of white and female students in less elite programs is reduced by [] percent. 

My results indicates that the removal of standardized exams shifts grades unevenly upwards, but the magnitude of the increments in grades do not necessary improve access to elite universities. The nonstandardized grading scheme rescales the preexisting measure of academic performance of students, with the mean of the grade distribution moving upward. The resulting grade scheme is top-coded, failing to signal student quality to universities, especially for students at the top of the distribution. In addition, the upward shift in the grade distribution is accompanied by a reshuffling of achievement levels on their final exams across demographic groups of students. Student demographic groups sort to different schools (e.g. private or public) and their subject choice (e.g. quantitative or nonquantitative), which all receives different extent of grade inflation leading to a overall reshuffling of top grades. My results indicate that private schools unilaterally inflated their students' grades, indicating an increase in the number of students educated in private schools at the top of the grade scale. Female and white students moved upwards then other student groups, due to choosing exam subjects with more teacher leniency on grading. 

My paper points to how integrating non-standardizing testing into university admission changes the composition of universities with long-lasting impact on the graduates labor market outcomes.